\section{Motivation}

\subsection{Threat Model}

\subsubsection{The ontology identification problem}
We can't yet reliably look inside a model and draw conclusions about its
goals or world model; interpretability has found a model's hidden
objective in a planted test~\cite{marks2025}.
This is a major limitation on almost all proposed alignment agendas,
because behavior alone cannot settle the question: a model can behave
differently when it believes it is being trained, and when that belief
reaches it through training documents rather than its prompt, the
difference can persist with no written reasoning to
inspect~\cite{greenblatt2024}. Only some models show such a
gap~\cite{sheshadri2025}, but that a good record under test need not
predict later behavior is an old argument~\cite[Ch.~8]{bostrom2014}.

\paragraph{Inner misalignment}
We don't currently know how to specify a robust goal for a model.
Because goal space is large, we should expect
goal misgeneralization as the default outcome.


\subsubsection{Limits of natural abstractions}
Networks trained on the same task can share an algorithm across
architectures while differing, even between random seeds, in which
features implement it: evidence for weak but not strong
universality~\cite{chughtai2023}. In language models, SAE features of a
Transformer and a Mamba trained on the same data match nearly as well as
those of two Transformer seeds~\cite{wang2024}.
This implies the existence of at least a weak version of
natural abstractions~\cite{lang2023},
which would be a very convenient shortcut for auditing a model.

Natural abstractions could be considered the set of concepts we could use to communicate with aliens. 
They should consistently emerge through unsupervised learning (in the limit given some constraints).
If any neural network is considered a fuzzy generalization of a finite state machine, 
natural abstractions are the latent states.

Unfortunately, the abstractions a model learns seem highly sensitive to choice of training data, and even to the random seed with the data fixed~\cite{chughtai2023}, and are thus not necessarily universal in the limit.
\footnote{The same applies to those that seem most natural to humans.}
Assuming that model capabilities will continue to primarily be driven by access to more data and more modalities, we cannot expect the abstractions learned by a weak model to remain robust in a more powerful model.

Even if we could fully characterise a model's state space in terms of natural abstractions,
that still doesn't mean we can reconstruct the latent FSM, 
as the transition rules can be arbitrarily complex.


\subsection{Theory of change}
Conditional on the underlying model of ontology being true, the
techniques proposed here would be a substantial interpretability
advance. The agenda also has a case for counterfactual impact: its
theoretical basis consists of nonobvious applications of obscure
techniques from niche fields of bioinformatics, and similar research
is unlikely to be done soon otherwise.

\subsubsection{What a solution would look like}
Mechinterp plausibly offers the best path to a robust alignment
solution despite criticisms.
The case rests on
\begin{enumerate}
    \item \textbf{decomposability} Mechinterp allows alignment to be
    broken into manageable subproblems.
    \item \textbf{parallelizability} Decomposable subproblems can be
    worked on independently.
    \item \textbf{robustness} Mechinterp offers direct causal
    explanations of model behavior.
\end{enumerate}
The most straightforward mechinterp alignment solutions are
\begin{enumerate}
    \item eliciting latent knowledge (ELK)
    Directly inferring a model's ontology.
    \label{it:elk}
    \item mechanistic anomaly detection (MAD)
    Detecting when a model has ``unusual'' thoughts.
    \label{it:mad}
\end{enumerate}
Both of these can be considered
ontology identification subproblems. MAD requires mapping between the model's activations and
the model's ontology. ELK requires mapping between the model's ontology
and human ontology. These are hard but not intractable problems.
Recently, developmental interpretability
(devinterp) techniques have enabled identification of ``developmental
stages'' corresponding to learning specific
tasks~\cite{hoogland2024developmental}. This is strong evidence that
learning consists primarily of the network
grokking~\cite{nanda2023progress} discrete concepts. While these results
are promising for MAD, locating these emergent capabilities in the
model's activations remains an open problem.

Interpretability glosses broadly as ``let humans understand AI''. Getting
AI to understand humans better scales badly: making the model smarter
leads to subtler and subtler issues with larger and larger effects.
Between the two sits scalable oversight, in which models help humans
supervise models: people working with an unreliable dialog assistant
already outperform both the assistant and themselves
alone~\cite{bowman2022}, and interpretability is one of the channels such
oversight can draw on. And what we're doing is not just
letting humans understand AI. Because an LLM reflects humanity from its
training data, high quality conceptual interpretability provides
machine-comprehensible mathematical representations of every concept
that humanity cares about. We work on interp for the usual
corrigibility reasons, and also because outer alignment needs a stable
place to attach values: a goal defined over one of a model's concepts
must keep pointing at that concept as the model's representations are
revised~\cite[Ch.~12]{bostrom2014}. The preliminary results measure a
small version of this problem: past the first layer, no head survives a
change of seed (Section~\ref{sec:prelim-seed}). This attacks both.

\subsubsection{Macro-scale strategy}
Broadly, the goals are to
\begin{enumerate}
  \item formalize the notion of a model's latent concept space
  \item distill a model's latent concept space
  \item enforce decomposability of latent concept space
  \item enforce robustness of existing concepts as the dimensionality of
  the latent concept space increases
\end{enumerate}
A prerequisite for all of these is the ability to express a
metaontology
which can identify latent concepts in activation space.


